% Part: normal-modal-logic
% Chapter: completeness
% Section: truth-lemma

\documentclass[../../../include/open-logic-section]{subfiles}

\begin{document}

\olfileid{nml}{com}{tru}

\olsection{सत्यता पूर्वप्रमेय}

कॅनॉनिकल प्रतिमान~$\mModel{M^\Sigma}$ ची व्याख्या
अशी केली आहे की $\mSat{M^\Sigma}{!A}[\Delta]$
असेल तेव्हा आणि तेव्हाच $!A \in \Delta$.
!!{propositional
variable}s साठी
$V^\Sigma$ ची व्याख्याच हे थेट दाखवते.
मात्र ही समतुल्यता सर्व !!{formula}s साठी
खरी आहे याची पडताळणी करावी लागेल.
हे आपण विगमनाने करू. विगमन-पायरीत
विधानीय कारक असलेल्या !!{formula}s साठी
समतुल्यता सिद्ध करावी लागते (येथे
\olref[ccs]{prop:ccs-properties} वापरतो),
तसेच मोडल कारक असलेल्या सूत्रांसाठीही
(येथे \olref[mod]{sec} मधील निष्कर्ष वापरतो).

\begin{prop}[सत्यता पूर्वप्रमेय]
  \ollabel{prop:truthlemma}
  प्रत्येक !!{formula}~$!A$ साठी
  $\mSat{M^\Sigma}{!A}[\Delta]$ असेल तेव्हा
  आणि तेव्हाच $!A \in \Delta$.
\end{prop}

\begin{proof}
  $!A$ च्या रचनेवर विगमन करू.
  \begin{enumerate}
    \tagitem{prvFalse}{\indcase{!A}{\lfalse}{$\mSat/{M^\Sigma}{\lfalse}[\Delta]$
        हे \olref[syn][trw]{defn:mmodels} नुसार, आणि
        $\lfalse \notin \Delta$ हे
        \olref[ccs]{prop:ccs-properties}\olref[ccs]{prop:ccs-lfalse}
        नुसार.}}{}

    \tagitem{prvTrue}{\indcase{!A}{\ltrue}{$\mSat{M^\Sigma}{\ltrue}[\Delta]$
        हे \olref[syn][trw]{defn:mmodels} नुसार, आणि
        $\ltrue \in \Delta$ हे
        \olref[ccs]{prop:ccs-properties}\olref[ccs]{prop:ccs-ltrue}
        नुसार.}}{}

    \item \indcase{!A}{p}{$\mSat{M^\Sigma}{p}[\Delta]$
      असेल तेव्हा आणि तेव्हाच $\Delta \in
      V^\Sigma(p)$, हे \olref[syn][trw]{defn:mmodels}
      नुसार. तसेच $\Delta
      \in V^\Sigma(p)$ असेल तेव्हा आणि तेव्हाच
      $p \in \Delta$, हे~$V^\Sigma$ च्या
      व्याख्येनुसार.}

    \tagitem{prvNot}{\indcase{!A}{\lnot !B}{%
        \iftag{probNot}{सराव.}{$\mSat{M^\Sigma}{\lnot !B}[\Delta]$
          असेल तेव्हा आणि तेव्हाच
          $\mSat/{M^\Sigma}{!B}[\Delta]$
          (\olref[syn][trw]{defn:mmodels} नुसार);
          असेल तेव्हा आणि तेव्हाच $!B \notin \Delta$
          (विगमन गृहीतकानुसार);
          असेल तेव्हा आणि तेव्हाच
          $\lnot !B \in \Delta$
          (\olref[ccs]{prop:ccs-properties}\olref[ccs]{prop:ccs-lnot}
          नुसार).}}}{}

    \tagitem{prvAnd}{\indcase{!A}{!B \land !C}{%
        \iftag{probAnd}{सराव.}{$\mSat{M^\Sigma}{!B \land
            !C}[\Delta]$ असेल तेव्हा आणि तेव्हाच
          $\mSat{M^\Sigma}{!B}[\Delta]$ आणि
          $\mSat{M^\Sigma}{!C}[\Delta]$
          (\olref[syn][trw]{defn:mmodels} नुसार);
          असेल तेव्हा आणि तेव्हाच $!B \in \Delta$
          आणि $!C \in \Delta$
          (विगमन गृहीतकानुसार);
          असेल तेव्हा आणि तेव्हाच $!B \land !C \in
          \Delta$
          (\olref[ccs]{prop:ccs-properties}\olref[ccs]{prop:ccs-land}
          नुसार).}}}{}

    \tagitem{prvOr}{\indcase{!A}{!B \lor !C}{%
        \iftag{probOr}{सराव.}{$\mSat{M^\Sigma}{!B \lor
            !C}[\Delta]$ असेल तेव्हा आणि तेव्हाच
          $\mSat{M^\Sigma}{!B}[\Delta]$ किंवा
          $\mSat{M^\Sigma}{!C}[\Delta]$
          (\olref[syn][trw]{defn:mmodels} नुसार);
          असेल तेव्हा आणि तेव्हाच $!B \in \Delta$
          किंवा $!C \in \Delta$
          (विगमन गृहीतकानुसार);
          असेल तेव्हा आणि तेव्हाच $!B \lor !C \in
          \Delta$
          (\olref[ccs]{prop:ccs-properties}\olref[ccs]{prop:ccs-lor}
          नुसार).}}}{}

    \tagitem{prvIf}{\indcase{!A}{!B \lif !C}{%
        \iftag{probIf}{सराव.}{$\mSat{M^\Sigma}{!B \lif
            !C}[\Delta]$ असेल तेव्हा आणि तेव्हाच
          $\mSat/{M^\Sigma}{!B}[\Delta]$ किंवा
          $\mSat{M^\Sigma}{!C}[\Delta]$
          (\olref[syn][trw]{defn:mmodels} नुसार);
          असेल तेव्हा आणि तेव्हाच $!B \notin \Delta$
          किंवा $!C \in \Delta$
          (विगमन गृहीतकानुसार);
          असेल तेव्हा आणि तेव्हाच $!B \lif !C
          \in \Delta$
          (\olref[ccs]{prop:ccs-properties}\olref[ccs]{prop:ccs-lif}
          नुसार).}}}{}

    \tagitem{prvIff}{\indcase{!A}{!B \liff !C}{%
        \iftag{probIff}{सराव.}{$\mSat{M^\Sigma}{!B \liff
            !C}[\Delta]$ असेल तेव्हा आणि तेव्हाच
          $\mSat{M^\Sigma}{!B}[\Delta]$ आणि
          $\mSat{M^\Sigma}{!C}[\Delta]$
          ही दोन्ही, किंवा
          $\mSat/{M^\Sigma}{!B}[\Delta]$ आणि
          $\mSat/{M^\Sigma}{!C}[\Delta]$
          ही दोन्ही
          (\olref[syn][trw]{defn:mmodels} नुसार);
          असेल तेव्हा आणि तेव्हाच $!B \in \Delta$
          आणि $!C \in \Delta$ ही दोन्ही, किंवा
          $!B \notin \Delta$ आणि $!C \notin
          \Delta$ ही दोन्ही
          (विगमन गृहीतकानुसार);
          असेल तेव्हा आणि तेव्हाच $!B \liff !C \in
          \Delta$
          (\olref[ccs]{prop:ccs-properties}\olref[ccs]{prop:ccs-liff}
          नुसार).}}}{}

    \tagitem{prvBox}{\indcase{!A}{\Box !B}{%
        \iftag{probBox}{सराव.}{प्रथम
          $\mSat{M^\Sigma}{\Box !B}[\Delta]$
          असे समजा.
          \olref[syn][trw]{defn:mmodels} नुसार
          $R^\Sigma \Delta\Delta'$ असलेल्या
          प्रत्येक $\Delta'$ साठी
          $\mSat{M^\Sigma}{!B}[\Delta']$.
          विगमन गृहीतकानुसार
          $R^\Sigma \Delta\Delta'$ असलेल्या
          प्रत्येक $\Delta'$ साठी $!B \in
          \Delta'$. $R^\Sigma$ च्या
          व्याख्येनुसार $\Box^{-1} \Delta
          \subseteq \Delta'$ असलेल्या प्रत्येक
          $\Delta'$ साठी $!B \in \Delta'$.
          \olref[mod]{prop:box} नुसार
          $\Box !B \in \Delta$.

          आता $\Box !B \in \Delta$ असे समजा.
          $R^\Sigma \Delta\Delta'$ म्हणजेच
          $\Box^{-1}\Delta \subseteq \Delta'$
          असलेला कोणताही $\Delta' \in W^\Sigma$
          घ्या. $\Box !B \in \Delta$
          असल्याने $!B \in \Box^{-1} \Delta$.
          परिणामी $!B \in \Delta'$.
          विगमन गृहीतकानुसार
          $\mSat{M^\Sigma}{!B}[\Delta']$.
          $R^\Sigma \Delta\Delta'$ असलेला
          $\Delta'$ स्वेच्छ असल्याने
          $R^\Sigma \Delta\Delta'$ असलेल्या
          सर्व $\Delta' \in W^\Sigma$ साठी
          $\mSat{M^\Sigma}{!B}[\Delta']$.
          \olref[syn][trw]{defn:mmodels} नुसार
          $\mSat{M^\Sigma}{\Box
            !B}[\Delta]$.}}}{}

    \tagitem{prvDiamond}{\indcase{!A}{\Diamond !B}{%
        \iftag{probDiamond}{सराव.}{प्रथम
          $\mSat{M^\Sigma}{\Diamond !B}[\Delta]$
          असे समजा.
          \olref[syn][trw]{defn:mmodels} नुसार
          $R^\Sigma \Delta\Delta'$ असलेल्या
          एखाद्या $\Delta'$ साठी
          $\mSat{M^\Sigma}{!B}[\Delta']$.
          विगमन गृहीतकानुसार
          $R^\Sigma \Delta\Delta'$ असलेल्या
          एखाद्या $\Delta'$ साठी
          $!B \in \Delta'$.
          $R^\Sigma$ च्या व्याख्येनुसार,
          \iftag{prvBox}{$\Box^{-1} \Delta
            \subseteq \Delta'$ असलेल्या
            एखाद्या $\Delta'$ साठी
            $!B \in \Delta'$.
            \olref[mod]{prop:diamond} नुसार,}{}
          $\Diamond \Delta' \subseteq \Delta$
          असलेल्या एखाद्या $\Delta'$ साठी
          $!B \in \Delta'$.
          $!B \in \Delta'$ असल्याने
          $\Diamond !B \in \Diamond
          \Delta'$; म्हणून
          $\Diamond !B \in \Delta$.

          आता $\Diamond !B \in \Delta$
          असे समजा. \olref[mod]{prop:diamond}
          नुसार $\Diamond \Delta' \subseteq
          \Delta$ आणि $!B \in \Delta'$
          असलेला संपूर्ण $\Sigma$-सुसंगत
          $\Delta' \in W^\Sigma$ असतो.
          \iftag{prvBox}{\olref[mod]{lem:box-iff-diamond}
            नुसार $\Box^{-1}
            \Delta \subseteq \Delta'$
            आणि $!B \in \Delta'$ असलेला
            $\Delta' \in W^\Sigma$ असतो.}{}
          $R^\Sigma$ च्या व्याख्येनुसार
          $R^\Sigma \Delta\Delta'$;
          म्हणून $R^\Sigma
          \Delta\Delta'$ आणि
          $!B \in \Delta'$ असलेला
          $\Delta' \in W^\Sigma$ असतो.
          \olref[syn][trw]{defn:mmodels}
          नुसार $\mSat{M^\Sigma}{\Diamond
            !B}[\Delta]$.}}}{}
  \end{enumerate}
\end{proof}

\begin{probtag}{probFalse,probTrue,probNot,probAnd,probOr,probIf,probIff,probBox,probDiamond}
  \olref[nml][com][tru]{prop:truthlemma}
  ची सिद्धता पूर्ण करा.
\end{probtag}

\end{document}
